\section{The odd-dimensional genuine Spin boundary}
\label{sec:spin-odd-extension}

\paragraph{Status and exact scope.}
This is a separate new author derivation, D-OL-22--24.
The complete root contributor reading and the separate
/root/g34\_eightfolds whole-proof/source/code audit have passed.
Root supplied the order-four strategy; independent origin is not claimed
for that contribution. Distinct exact carrier reconstructions are retained.  The field is odd,
the natural dimension is odd, and the module is any nontrivial
irreducible complex module of the finite Spin group, with an
arbitrary central kernel.  Descent to the natural $\Omega$ group,
extension to a full orthogonal group, selfduality and a determinant
condition are not assumed.  Earlier canonical statements are unchanged.

Let $F=\mathbf F_q$, $q$ odd, and let $(V,Q)$ be nondegenerate.
Put $S=\operatorname{Spin}(V)$, and denote its natural map by
$\varphi:S\longrightarrow\Omega(V)$.  We use
\[
 \beta=\frac{4997}{30000},\qquad
 \delta=\frac{1-(19/20)^9}{2},\qquad
 \frac16<\delta<\frac3{16}<\frac7{32},
 \qquad \nu(T)=a-\max_\lambda\dim\ker(T-\lambda I).
 \tag{OSP.1}
\]

\paragraph{The precise lift and radical interface.}
We use D-OL-19 and the parity-free opening arguments of D-OL-20.
Their full proofs give the following interface in either parity
for $d=\dim V\ge22$.  Every compatible normalizer of a represented
central quotient of $S$ is implemented by a natural semilinear
similitude $H$, and the two projective orders agree.
For clarity the lift on even Clifford words is
\[
 \Phi_H(v_1\cdots v_{2k})
       =\mu^{-k}Hv_1\cdots Hv_{2k},
 \qquad Q(Hv)=\mu\,\sigma Q(v).                              \tag{OSP.2}
\]
The two possible tensor lengths in a Clifford relation differ by
two, and the extra $\mu^{-1}$ cancels the multiplier of $Q$.
Thus this formula is well defined on the even Clifford algebra.
The corresponding odd-word map, scaled by $\mu^{-k}$ on length
$2k+1$, intertwines its two even-module actions and restricts to
$H$ on $V$.  Reversal commutes with (OSP.2), so the Spin norm
and vector covariance are preserved.  This is the explicit
semilinear Clifford construction of D-OL-19; its even-algebra
calculation is the published construction
\cite[Part II, pp.~48--49]{SF:Wonenburger1962}.
No theorem asserting an unrestricted extension to the whole
Clifford algebra is used.

We recall the represented-kernel point.  The published exact sequence
and quasisimplicity input are
\[
 1\longrightarrow\{\pm1\}\longrightarrow S
       \mathop{\longrightarrow}^{\varphi}\Omega(V)
       \longrightarrow1,                                  \tag{OSP.3}
\]
with $S$ perfect in these dimensions
\cite[proof of Theorem~4.3.6, Step~1, pp.~1918--1919]{HOS:LST2011}.
If two homomorphisms from $S$ to a represented quotient induce the
same simple-quotient automorphism, their ratio is a homomorphism
to the center; perfectness makes it trivial.  Applying this to
the actual normalizing operator and (OSP.2) proves preservation
of its actual central kernel, rather than assuming it.
The full natural semilinear image, including the scalar
centralizer and arbitrary field components, is D-OL-07 and its
complete outer-order comparison
\cite[Section~2.2.2, Table~2.1, p.~14]{HES:LiXia2022}.

Every odd-order one- or two-plane radical $U\subset\Omega(V)$
has a unique lift $\widetilde U\subset S$ isomorphic to it.
Indeed a Sylow $p$ subgroup of $\varphi^{-1}(U)$ maps
isomorphically to $U$, and the preimage is its product with
the central kernel $\{\pm1\}$.  That Sylow subgroup is
therefore normal and unique.  All natural conjugation formulas
on $U$ consequently lift exactly, without a central phase.
If opposite radicals generate $\Omega(V)$, their lifts
generate $S$: their generated subgroup maps onto $\Omega(V)$,
so its quotient in $S$ is central and abelian, and perfectness
eliminates the quotient.  Opposite conjugate radicals therefore
have equal fixed-space dimensions with zero intersection on a
nontrivial irreducible $S$-module.  Each fixed dimension is at
most $a/2$.  Using the full inverse image of a natural Levi
group gives all its natural weight orbits and the constancy of
weight multiplicities.  This does not assert that the module
factors through that Levi quotient.

\subsection*{D-OL-22. The fourth-order negative-reflection carrier}

\begin{lemma}[D-OL-22]
Let $q=3$, let $d\ge25$ be odd, and suppose
$g\in SO(V)$ satisfies $g^2=-r_v$, where $v$ is nonsingular
and $r_v$ is its reflection.  If a finite-projective-order
normalizer $T$ of a nontrivial irreducible $S$-module implements
this natural conjugation, then
\[
                              \nu(T)\ge7a/32.               \tag{OSP.4}
\]
The assertion includes $g\notin\Omega(V)$.
\end{lemma}

\begin{proof}
The element $g$ preserves $Fv$ and $W_0=v^\perp$.
On $W_0$ it satisfies $g^2=-I$.  Since $x^2+1$ is irreducible
over $\mathbf F_3$, its characteristic polynomial there is
$(x^2+1)^{(d-1)/2}$, with determinant one.
The determinant-one condition on $g$ therefore makes its
eigenvalue on $Fv$ equal to one.

Put $E=F[t]/(t^2+1)$ and make $W_0$ an $E$-space by
letting $t$ act as $g$.  If $B$ is the polar form, then
$g^*=-g$ on $W_0$, and
\[
 b(x,y)=\frac12\bigl(B(x,y)+t^{-1}B(gx,y)\bigr)              \tag{OSP.5}
\]
is a nondegenerate $E/F$-Hermitian form with trace $B$.
Multiplication by $t$ in its first argument multiplies it by
$t$, and in its second by $-t$; symmetry of $B$ gives the
Hermitian identity.  The finite-field Hermitian hyperbolic-pair
theorem applies since $\dim_EW_0=(d-1)/2\ge12$
\cite[Lemma~10.2, Corollary~10.3 and ensuing decomposition,
pp.~116--117]{OS:Taylor1992}.
Its opposite $E$-lines give opposite $g$-stable totally
singular $F$-planes $L,L'$.

Choose dual bases and write
\[
 \begin{gathered}
 V=L\oplus W\oplus L',\qquad N=\dim W=d-4\ge21,\\
 Q(x,w,y)=x^{\mathsf t}y+\tfrac12w^{\mathsf t}Bw,\qquad
 g=\operatorname{diag}(A,h,A^{-\mathsf t}),\\
 A^2=-I_2,\quad \det A=1,\quad
 W=Fv\perp W_1,\quad h|_{Fv}=1,\quad h^2|_{W_1}=-I.
 \end{gathered}                                             \tag{OSP.6}
\]
The complementary dimension $N$ is odd; no even-dimensional
assumption on $W$ is being imported.
The natural two-plane radical consists of
\[
 u(X,Z)=
 \begin{pmatrix}
 I_2&X&Z-\frac12XB^{-1}X^{\mathsf t}\\
 0&I_N&-B^{-1}X^{\mathsf t}\\
 0&0&I_2
 \end{pmatrix},
 \quad X\in M_{2,N}(F),\quad Z^{\mathsf t}=-Z.               \tag{OSP.7}
\]
Its group law has central correction
$\frac12(YB^{-1}X^{\mathsf t}-XB^{-1}Y^{\mathsf t})$.
This is the same Heisenberg group of order $q^{2N+1}$
as in D-OL-16.  The two opposite such radicals generate
$\Omega(V)$ by its complete line-radical generation argument,
which uses only a sufficiently large nondegenerate complement
and applies in either parity.  The preceding lift interface
therefore gives
\[
                     t_0:=\dim M^{\widetilde Q_L}\le a/2.   \tag{OSP.8}
\]

Conjugation by $g$ on (OSP.7) is exactly
\[
 \mathcal A:X\longmapsto AXh^{-1},\qquad
 Z\longmapsto AZA^{\mathsf t}=Z.                            \tag{OSP.9}
\]
The center is fixed.  Decompose
$\operatorname{Hom}_F(W,L)$ according to $W=Fv\oplus W_1$.
On the first summand, of dimension two, the action is $A$,
with square $-I$ and no fixed vector.  On the second summand
its square is the identity.  The fixed maps there are
$E$-linear maps $W_1\to L$, of $F$-dimension $N-1$.
It follows that the exact fixed dimensions of
$\mathcal A,\mathcal A^2,\mathcal A^3$ on the $2N$-dimensional
quotient are
\[
                              N-1,\quad 2N-2,\quad N-1.    \tag{OSP.10}
\]

The dual weight action is $Y\mapsto hYA^{-1}$ on
$\operatorname{Hom}_F(L,W)$.  A fixed map has image in $W_1$
and is $E$-linear, so the fixed space has dimension $N-1$.
Every nonzero fixed map has rank two: its kernel would otherwise
be a proper nonzero invariant subspace of the irreducible
two-dimensional $F[A]$-module $L$.

The rank-two Levi orbits are the ordered independent pairs
with prescribed Gram matrix, by Witt extension and the
determinant/spinor corrections in their large perpendicular
space.  The complete estimate of D-OL-16 applies for odd $N$
as well.  Explicitly, the first-vector and affine second-vector
counts, after excluding dependent vectors, are bounded below by
\[
 q^{N-1}-q^{\lceil N/2\rceil}-1,\qquad
 q^{N-2}-q^{N/2}-q.                                       \tag{OSP.11}
\]
The affine count uses polar rank at least $N-2$; if a
linear term is nonzero on the radical, its Gauss sum is zero.
Thus it also covers degenerate prescribed Gram matrices.
Each bound exceeds half its leading term for $q=3$, $N\ge21$.
Every rank-two orbit consequently has more than $q^{2N-3}/4$
vectors, whereas there are at most $q^{N-1}$ fixed dual maps.
Its fixed proportion is less than
\[
                              4q^{-N+2}<1/8.               \tag{OSP.12}
\]
Rank-one orbits have no fixed map.  Weight multiplicities are
constant on these full-preimage Levi orbits.
The action in (OSP.9) has order four, so every nonfixed
weight orbit has length two or four.  Each such orbit contributes
at most half its isotypic dimension to a specified eigenspace.
Thus on the sum of all nontrivial linear-character spaces of
$\widetilde Q_L$ the support is at least $7/16$ of that sum.
It is not necessary that an individual Levi orbit be invariant
under $\mathcal A$; the fixed-weight bound and weight
multiplicities give the same bound after summing all of them.

It remains to retain the nonlinear Heisenberg sectors and
their multiplicities.  For each nontrivial additive central
character $\psi$, D-HOS-01 gives its unique carrier $U_\psi$,
of degree $q^N$, with Weyl operators $L_X$ forming a basis
of $\operatorname{End}(U_\psi)$.  The center is fixed, so
choose a unitary carrier intertwiner $\Xi$ and rescale it to
$\Xi^4=I$.  The exact zero-central-coordinate action (OSP.9)
gives, for $j=1,2,3$,
\[
 \Xi^jL_X\Xi^{-j}=L_{\mathcal A^jX},\qquad
 |\operatorname{tr}\Xi^j|^2
   =\operatorname{tr}(\operatorname{Ad}\Xi^j)
   =|\ker(\mathcal A^j-I)|.                                \tag{OSP.13}
\]
There is no hidden central phase, by uniqueness of the
odd-order radical lift.  Equations (OSP.10)--(OSP.13) give
the trace magnitudes
$q^{(N-1)/2},q^{N-1},q^{(N-1)/2}$.
The four-term Fourier formula therefore bounds every
eigenspace of $\Xi$ by
\[
 \frac{q^N+q^{N-1}+2q^{(N-1)/2}}4
       <\frac9{16}q^N.                                    \tag{OSP.14}
\]
For $q=3$, $N\ge21$ the strict inequality follows directly
after division by $q^N$; the last small term is at most
$2\cdot3^{-11}$.  Thus its support fraction exceeds $7/16$.
On the full $\psi$-isotypic space, the actual operator is
$\Xi\otimes C$ by the matrix-algebra commutant.  Rescaling the
given finite-projective-order $T$ to finite order makes the
multiplicity operator diagonalizable.  Eigenvalue convolution
preserves the carrier support bound for every scalar phase.

Only the trivial-radical space of dimension $t_0$ remains
uncontrolled.  Combining all sectors with (OSP.8) proves
$\nu(T)\ge(7/16)(a-t_0)\ge7a/32$.  This proves the lemma.
\end{proof}

\subsection*{D-OL-23. All projective orders and all compatible normalizers}

\begin{theorem}[D-OL-23]
Let $q$ be odd, let $d\ge1135$ be odd, and let $M$ be any
nontrivial irreducible complex $S=\operatorname{Spin}_d(q)$-module
of degree $a$.  Every nonscalar compatible finite-projective-order
normalizer $T$ satisfies $\nu(T)\ge\beta a$.
For $q\ge5$ it satisfies $\nu(T)>\delta a$.
For $q=3$, an operator with $\nu(T)\le\delta a$ has natural
projective lift a negative reflection, and projective order two.
Moreover
\[
                   a\ge q^{d-3}-q^{\lceil(d-2)/2\rceil}-1>360.
 \tag{OSP.15}
\]
\end{theorem}

\begin{proof}
The degree estimate is the parity-free argument of D-OL-20,
which we recall to avoid an ordinary-character hypothesis.
Some nontrivial weight of a lifted full line radical occurs;
otherwise its conjugate opposite would also act trivially,
and their generation of $S$ would make $M$ trivial.
The full inverse-image Levi moves that weight through a
nonzero natural norm orbit of dimension $d-2$.  Its size is
at least the first expression in (OSP.15), by the same exact
quadratic Gauss count as in D-OL-06.  Every weight in the orbit
occurs with positive multiplicity.  This proves (OSP.15) for
every Spin character, without applying an ordinary $\Omega$
degree theorem.

Take a nonscalar power of $T$ having prime projective order
$\ell$; this cannot increase $\nu$.  The natural lift supplied
by D-OL-19 has the same projective order.  If its field component
is nontrivial, $q=q_0^\ell$.  The complete norm and quadratic-form
descent (ON.7)--(ON.9) rescales this possibly mixed linear-field
lift to order $\ell$ with an invariant hyperbolic pair over
$\mathbf F_{q_0}$.  Its action on the associated line-radical
weights is Frobenius, with fixed proportion less than $1/4$.
The unique lifted radical and full inverse-image Levi retain
this action exactly.  The resulting support is at least
$3a/16>\delta a$.  No pure-field assumption is made.

Suppose the prime-projective-order lift is linear.
In odd natural dimension a similitude multiplier is a square,
since $(\det H)^2=\mu^d$.  Rescale it to an isometry and,
if necessary, multiply by $-I$ to obtain the unique
representative $g\in SO(V)$.  The only natural scalar in
$SO(V)$ is $I$.  Thus the order of $g$ is exactly the
projective order of $T$.

If $\ell$ is odd, $g$ has trivial spinor norm and belongs to
$\Omega(V)$, since $SO(V)/\Omega(V)$ has order two.
Choose a Spin lift $\widetilde g$.  Its conjugation on the
full Clifford algebra implements $g$ on $V$, hence agrees
on the even algebra with (OSP.2).  Schur's lemma makes the
normalizing operator a scalar multiple of $\rho(\widetilde g)$.
If $g$ fixes a hyperbolic pair, the complementary natural
moving rank is at least two: a rank-one nontrivial isometry
is a determinant-minus-one reflection.  The all-rank line
Gauss count (ON.3), with the lifted radical, gives $7a/32$.

If no hyperbolic pair is fixed, the natural support restricted
to $\{1,-1\}$ is at least $567$.  For semisimple odd prime
order, the fixed space is nondegenerate of dimension at most
two and $-1$ is not an eigenvalue.  For characteristic prime
order, if $r=\operatorname{rank}(g-I)$, the fixed polar rank
is at least $d-2r$.  Absence of a hyperbolic pair forces it
to be at most two, so the integer $r\ge(d-2)/2$ is at least
$567$.  Again $-1$ is not an eigenvalue.
Apply the published estimate directly to the nontrivial
irreducible Spin character:
\[
 |\chi(\widetilde g)|/a<(19/20)^9,\qquad
                     \nu(\rho(\widetilde g))>\delta a.      \tag{OSP.16}
\]
The exact input is
\cite[Theorem~4.3.6, (4.3.2), pp.~1917--1918;
complete proof pp.~1918--1923; Remark~4.3.7, p.~1923]{HOS:LST2011}.
Its restricted scalar set for Spin is $\mu_2$ and its
threshold is $567=7\cdot9^2$.  A central choice of lift changes
neither the natural action nor the absolute character bound.
The passage to support uses
$|\chi|\ge a-2\nu$, after rotating a largest-eigenspace
eigenvalue to one.  No faithful-character hypothesis or
descent to $\Omega$ is imposed.

If $\ell=2$, the normalized $g\in SO(V)$ is an involution.
Its two eigenspaces are nondegenerate and at least one
contains a hyperbolic pair.  Choose such a pair in the
$\varepsilon$-eigenspace.  The line-radical action on its
complement $W$ is the isometry $A=\varepsilon g|_W$,
exactly as in (OO.5), and it is nonidentity.
For moving rank at least two the bound is $7a/32$.
Moving rank one means that $A$ is a reflection and
$g$ is projectively a reflection.  The unique $SO(V)$
representative in odd dimension is $g=-r_v$.
The precise rank-one Gauss fraction is bounded by
\[
 c_q=\frac{q^{-1}+q^{-9}}{1-q^{-9}-q^{-19}}.
 \tag{OSP.17}
\]
For $q\ge5$, $c_q<1/4$; for $q=3$, $c_3<1-4\beta$.
The lifted-radical fixed-space bound and the two-element
weight orbits give support at least $(1-c_q)a/4$.
This is greater than $3a/16$ in the first case and
greater than $\beta a$ in the second.
These arguments treat every prime projective order and
give the claimed universal bounds.

It remains to classify every low-support order when $q=3$.
There is no field automorphism.  Let $g\in SO(V)$ be the
unique natural representative of the original projective
normalizer; its order equals that projective order.
An odd prime divisor would yield a strong-bound power by
the preceding analysis, so the order is $2^s$.
Its involution power must be $-r_v$ if the support is at
most $\delta a$.  If $s=1$, the claimed conclusion holds.

Suppose $s\ge3$.  Then $g^2\in\Omega(V)$ and
\[
                 (g^2)^{2^{s-2}}=-r_v .
 \tag{OSP.18}
\]
The exponent is even.  Any vector with eigenvalue $1$
or $-1$ for $g^2$ is therefore fixed by $-r_v$, hence
lies in $Fv$.  Both restricted eigenspaces have dimension
at most one.  The natural $\mu_2$-support of $g^2$ is
at least $d-1\ge567$.  The operator $T^2$ is a scalar
multiple of a Spin operator above $g^2$, by the exact
inner compatibility already used.  The same published
bound (OSP.16) applies to this element without a
prime-order condition.  It gives
$\nu(T)\ge\nu(T^2)>\delta a$, a contradiction.

The remaining case is $s=2$, so $g^2=-r_v$.
D-OL-22 applies whether $g$ belongs to $\Omega(V)$
or its other $SO$ coset, and gives
$\nu(T)\ge7a/32>\delta a$, another contradiction.
Thus $s=1$ and the natural projective lift is a negative
reflection.  The actual complex operator may have a
nontrivial central scalar in its square; its \emph{projective}
order, and its two normalized eigenspaces, are exactly two.
This distinction is retained in the source argument below.
\end{proof}

\subsection*{D-OL-24. The odd-dimensional Spin source contradiction}

\begin{theorem}[D-OL-24]
Allow a nonempty irreducibly acting normal central product of
represented Spin factors of odd natural dimensions $d_i\ge1135$
over arbitrary odd fields.  Give each factor any nontrivial
irreducible complex module, with arbitrary central kernel.
All compatible scalars, outer automorphisms and factor
permutations are allowed.  Such an actual inducing image
cannot supply period-image rank nine for its six-dimensional
multiplicity variation.  Transitivity, selfduality and a
determinant hypothesis are not required for this exclusion.
\end{theorem}

\begin{proof}
Write $e=\prod_i a_i$.  The degree bound (OSP.15)
gives $e>360$.  Convolution retains the support fraction
of every active factor-preserving operator.
Nontrivial factor cycles have projective support at least
$e/4$ by the complete weighted-word argument used in
D-OL-18.  A nonidentity scalar has actual moving rank $e$.
Every genuine inertia therefore moves at least $\beta e$.

Call an inertia good if its actual moving rank is at
least $\delta e$ and bad otherwise.  A bad operator is
factor preserving.  Every active factor must have $q_i=3$
and natural projective lift a negative reflection, by
D-OL-23.  Normalize each active factor's scalar so that
its two eigenvalues are $1,-1$.  Their proportions lie
in $[\beta,1-\beta]$, regardless of the original Spin
center or of the unnormalized square of the operator.
For two active factors the largest tensor-eigenspace
proportion is at most $(1-\beta)^2+\beta^2$:
both possible convolution expressions are affine in each
proportion and are maximized at a corner of that rectangle.
Thus two active factors already have support at least
\[
                   2\beta(1-\beta)e>e/4>\delta e.            \tag{OSP.19}
\]
Additional factors cannot increase the largest
eigenvalue proportion.  Every bad inertia consequently
has exactly one active factor.  Its actual scalar has to
select a fixed eigenvalue; otherwise it moves all $e$
dimensions.  No assertion about scalar preservation of
a nonreal-eigenvalue estimate is used.

The actual intermediate-cover Euler identity is
$6=2e(h-1)+D$.  Since $e>360$, it excludes $h\ge2$.
For $h=1$ there are at most six genuine marks and the
finite-cover period image has rank at most six.
The proposed rank-nine case must therefore have a sphere
base, at least twelve genuine marks, a generating
product-one inertia tuple and $D=2e+6$.
The exact inequalities
\[
 13\beta-2>\frac17,\qquad
 11\beta+\delta-2>\frac1{60},\qquad
 2\beta(1-\beta)>\frac14                                  \tag{OSP.20}
\]
exclude thirteen marks and exclude even one good inertia
among twelve.  Indeed they would force respectively
$e<42$ or $e<360$.  Thus exactly twelve bad inertias remain.

These actual inertia elements generate the inducing image.
All preserve every factor.  Fix one normal factor; the
induced projection contains its full simple inner group
$\Omega(V_i)$, because the represented Spin factor is
contained in the inducing group and
$\Omega(V_i)$ has trivial natural center in odd dimension.
If $q_i\ne3$, all bad inertias have trivial projection
there, an immediate contradiction.  For $q_i=3$ their
projections are identities or projective reflections.
Choose at most twelve natural reflections with
nonsingular normals $v_j$ representing those projections.
Their common fixed subspace
\[
          U_0=\bigcap_jv_j^\perp,\qquad
                         \dim U_0\ge d_i-12>0              \tag{OSP.21}
\]
is fixed by every word in them.  Every element of the
natural $\Omega(V_i)$ is projectively such a word,
so it acts by an original-field scalar on $U_0$.
Those scalars form a homomorphism from the perfect group
$\Omega(V_i)$ to $F_3^*$.  They must all be one.
This would make $U_0$ a nonzero natural fixed space of
$\Omega(V_i)$, contradicting its opposite full line
radicals and their generation.  This argument uses only
the natural simple projection; no splitting of a
represented Spin center has been introduced.
The theorem follows.
\end{proof}

\paragraph{Consequences and remaining scope.}
Together with D-OL-21, the same common support and
reflection-generation argument also excludes mixed
odd-characteristic Spin factors of both parities when
each odd dimension is at least $1135$ and each even
dimension is at least $1136$.
In the transitive-core main-source problem this gives
the necessary bound $d\le1133$ in odd natural dimension
and $d\le1134$ in even natural dimension, without an
ordinary-module restriction.  These are necessary
parameter bounds and imply no existence at the
remaining dimensions.

The new argument is the fourth-order carrier and its
all-order source consequence.  The exact finite
supplement verifies natural matrices, all three fixed
dimensions, finite Weyl covariance, carrier trace
counts and the rational constants.  It does not
replace the universal Clifford lift, group generation,
finite-field Witt theory, published character estimate
or actual-cover argument.  The primary lifting and
quasisimplicity interfaces are the fully written
D-OL-19--20 proofs, with their precise source records.
No additional Schur-multiplier classification, full
orthogonal reality theorem or positive cycle supplier
is asserted here.
